\section{Implementation and Evaluation}
\label{sec:implementation}

We implement \(\protocol\) and evaluate its performance. The implementation follows the construction in Section~\ref{sec:construction}.  We instantiate the CP-SNARK with Groth16, using the commit-and-prove
interface~\cite{gnark,legosnark}.  For CP-Link, we use a pairing-based QA-NIZK
proof~\cite{qa-nizk}; the concrete instantiation used in our implementation is
given in Appendix~\ref{app:cplink-instant}.  The implementation is available at
\url{https://anonymous.4open.science/r/LAMP-38E3/.}

Our evaluation is organized around four questions.  First, does the number of  constraints of \(\protocol\) scale linearly with the matrix dimension, as predicted by Section~\ref{sec:protocol-complexity}?  Second, which components dominate proving time, proof size, and verification time?  Third, how does \(\protocol\) behave when proving multiple independent claims \(\mat{A}_i\mat{B}_i=\mat{C}_i\) in a batch, rather than a single product \(\mat{A}\mat{B}=\mat{C}\)?  Fourth, does the advantage remain useful on a neural-network workload? \\

\noindent\textbf{Experimental setup.}
We implement both \(\protocol\) and the Freivalds baseline in Go using gnark, with BN254 as the elliptic curve.  
All experiments were run on a Linux server with an AMD EPYC 7B13 CPU with 32 cores and 256GB of memory.
\rev{Each experiment was repeated ten times, and we report the average over these runs.}
We use a Reed--Solomon code of rate \(\rho=1/2\) and set the number of sampled positions to \(t=\revmath{309}\).\




\begin{table*}[!t]
\centering
  \ifshowrevisions
    \rowcolors{1}{yellow!20}{yellow!20}
  \fi
\setlength{\tabcolsep}{3.5pt}
\begin{adjustbox}{max width=\textwidth}
\begin{tabular}{clrrrr}
\toprule
\(\boldsymbol{k}\) & \textbf{System} & \textbf{Constraints}
& \textbf{Prove Time} & \textbf{Verify Time} & \textbf{Proof Size} \\
\midrule
           & \(\protocol\) & -- & -- & -- & -- \\
           & Freivalds & 49{,}610 & 0.59\,s & 2\,ms & 196\,B \\
\multirow{-3}{*}{$2^7$} & DualMatrix & -- & 3.57\,s & 0.13\,s & 71{,}288\,B \\
\midrule
           & \(\protocol\) & 329{,}377 & 2.94\,s & 0.13\,s & 50{,}308\,B \\
           & Freivalds & 197{,}194 & 2.02\,s & 2\,ms & 196\,B \\
\multirow{-3}{*}{$2^8$} & DualMatrix & -- & 4.81\,s & 0.14\,s & 77{,}140\,B \\
\midrule
           & \(\protocol\) & 655{,}245 & 5.73\,s & 0.14\,s & 65{,}860\,B \\
           & Freivalds & 788{,}810 & 6.99\,s & 2\,ms & 196\,B \\
\multirow{-3}{*}{$2^9$} & DualMatrix & -- & 7.98\,s & 0.15\,s & 82{,}992\,B \\
\midrule
           & \(\protocol\) & 1{,}306{,}972 & 12.56\,s & 0.13\,s & 80{,}772\,B \\
           & Freivalds & 3{,}156{,}298 & 25.91\,s & 1\,ms & 196\,B \\
\multirow{-3}{*}{$2^{10}$} & DualMatrix & -- & 18.54\,s & 0.15\,s & 88{,}844\,B \\
\midrule
           & \(\protocol\) & 2{,}609{,}193 & 27.49\,s & 0.13\,s & 99{,}012\,B \\
           & Freivalds & 12{,}622{,}154 & 102.32\,s & 2\,ms & 196\,B \\
\multirow{-3}{*}{$2^{11}$} & DualMatrix & -- & 57.02\,s & 0.16\,s & 94{,}696\,B \\
\midrule
           & \(\protocol\) & 5{,}214{,}877 & 69.95\,s & 0.13\,s & 119{,}556\,B \\
           & Freivalds & 50{,}495{,}818 & 411.28\,s & 2\,ms & 196\,B \\
\multirow{-3}{*}{$2^{12}$} & DualMatrix & -- & 205.21\,s & 0.17\,s & 100{,}548\,B \\
\midrule
           & \(\protocol\) & 10{,}426{,}236 & 193.89\,s & 0.13\,s & 136{,}260\,B \\
           & Freivalds & -- & -- & -- & -- \\
\multirow{-3}{*}{$2^{13}$} & DualMatrix & -- & 707.43\,s & 0.18\,s & 106{,}400\,B \\
\bottomrule
\end{tabular}
\end{adjustbox}
\vspace{6pt}
\caption{Performance comparison of \(\protocol\), the Freivalds baseline, and DualMatrix for square matrix multiplication.}
\label{tab:matmul-comparison}
\end{table*}


\subsection[Comparison with Matrix-Multiplication Proof Systems]
  {\rev{Comparison with Matrix-Multiplication Proof Systems}}
  \label{sec:eval-matmul-comparison}

  \begin{revblock}
  We compare \(\protocol\) with existing proof systems for matrix
  multiplication. Although zkMatrix and
  zkMaP are included in the theoretical comparison in
  Table~\ref{tab:asymptotic_comparison}, we were unable to identify publicly
  available implementations that would allow us to reproduce their experiments.
  We therefore compare \(\protocol\) against two systems: a Freivalds-based
  baseline implemented as an arithmetic circuit using
  LegoSNARK, and DualMatrix, for which
  a public implementation is available. Table~\ref{tab:matmul-comparison}
  reports the constraint count, proving time, verification time, and proof size
  of these systems. For both \(\protocol\) and DualMatrix, proving time includes
  commitment generation, and proof size includes the corresponding commitment
  data.

  For the \(\protocol\) experiments, we use a code rate of \(\rho=1/2\), which
  requires \(t=309\) queried positions. Since the codeword length is
  \(n=2k\), at \(k=2^7\) we have \(n=256<t\). Consequently, the
  \(\protocol\) experiments begin at \(k=2^8\). DualMatrix, on the other hand,
  is a specialized matrix-multiplication proof system that does not use an
  arithmetic circuit. An R1CS constraint count is therefore not an applicable
  metric for DualMatrix, and its constraint entries are left unspecified in the
  table.

  For the Freivalds baseline, the number of constraints increases by
  approximately \(4\times\) whenever the matrix dimension \(k\) doubles. This is
  because the vector--matrix products performed inside the circuit require
  \(\mathcal{O}(k^2)\) constraints. Moreover, its proving complexity is
  approximately \(\mathcal{O}(k^2\log k)\), causing the proving-time gap between
  the baseline and \(\protocol\) to grow with \(k\). At \(k=2^8\), the Freivalds
  baseline is faster, requiring \(2.02\) seconds compared with \(2.94\) seconds
  for \(\protocol\). Starting at \(k=2^9\), however, \(\protocol\) becomes
  faster. At \(k=2^{12}\), \(\protocol\) reduces the constraint count from
  \(50{,}495{,}818\) to \(5{,}214{,}877\), a \(9.68\times\) reduction, and
  reduces proving time from \(411.28\) seconds to \(69.95\) seconds, corresponding
  to a \(5.88\times\) speedup.

  The Freivalds baseline cannot complete the experiment at \(k=2^{13}\) because
  of insufficient memory. In contrast, \(\protocol\) completes the same
  experiment using approximately \(50.57\) GB of memory, with a proving time of
  \(193.89\) seconds. Because the Freivalds baseline uses the constant-size
  Groth16 proof provided by LegoSNARK, its proof size remains \(196\) B and its
  verification time remains approximately \(1\)--\(2\) ms, independently of the
  constraint count. Thus, the Freivalds baseline provides very small proofs and
  fast verification, but its proving time and memory consumption increase
  rapidly for large matrices.

  Across the range shared by \(\protocol\) and DualMatrix,
  \(2^8\leq k\leq2^{13}\), \(\protocol\) achieves a shorter proving time at every
  measured matrix dimension. At \(k=2^8\), \(\protocol\) requires \(2.94\)
  seconds, whereas DualMatrix requires \(4.81\) seconds, making \(\protocol\)
  approximately \(1.64\times\) faster. The gap increases with the matrix
  dimension: at \(k=2^{13}\), the respective proving times are \(193.89\) seconds
  and \(707.43\) seconds, giving \(\protocol\) a \(3.65\times\) speedup.
  Verification times are comparable and remain below \(0.2\) seconds for both
  systems, although \(\protocol\) is consistently slightly faster across all
  shared measurements.

  Proof size exhibits the opposite trend. From \(k=2^8\) through \(k=2^{10}\),
  \(\protocol\) produces smaller proofs than DualMatrix. Starting at
  \(k=2^{11}\), however, the proofs generated by \(\protocol\) become larger.
  For example, at \(k=2^{13}\), the proof sizes are \(136{,}260\) B for
  \(\protocol\) and \(106{,}400\) B for DualMatrix. These results demonstrate
  that \(\protocol\) provides faster proof generation and slightly faster
  verification than DualMatrix, at the cost of larger proofs for sufficiently
  large matrices.
  \end{revblock}


\begin{table*}[!t]
\centering
  \ifshowrevisions
    \rowcolors{1}{yellow!20}{yellow!20}
  \fi
% \scriptsize
\setlength{\tabcolsep}{3.5pt}
\begin{adjustbox}{max width=\textwidth}
\begin{tabular}{c cc ccc ccc ccc ccc}
\toprule
&
\multicolumn{2}{c}{\textbf{Setup Time}}
& \multicolumn{3}{c}{\textbf{Commit Time}}
& \multicolumn{3}{c}{\textbf{Prove Time}}
& \multicolumn{3}{c}{\textbf{Verify Time}}
& \multicolumn{3}{c}{\textbf{Proof Size}} \\
\cmidrule(lr){2-3}\cmidrule(lr){4-6}\cmidrule(lr){7-9}\cmidrule(lr){10-12}\cmidrule(lr){13-15}
\multirow{-2}{*}{\(\boldsymbol{k}\)}
& Groth16 & CP-Link
& Encode & ABC & XYZ
& Merkle & Groth16 & CP-Link
& Groth16 & Merkle & CP-Link
& Merkle & Groth16 & CP-Link \\
\midrule
$2^8$    & 20.03\,s  & 1.28\,s  & 0.04\,s  & 0.21\,s & 0.02\,s & 1\,ms & 2.59\,s  & 0.08\,s & 2\,ms & 1\,ms & 0.13\,s & 49{,}920\,B  & 260\,B & 128\,B \\
$2^9$    & 39.96\,s  & 2.52\,s  & 0.14\,s  & 0.74\,s & 0.05\,s & 1\,ms & 4.68\,s  & 0.12\,s & 2\,ms & 2\,ms & 0.14\,s & 65{,}472\,B  & 260\,B & 128\,B \\
$2^{10}$ & 80.58\,s  & 5.01\,s  & 0.46\,s  & 2.66\,s & 0.14\,s & 1\,ms & 9.09\,s  & 0.21\,s & 2\,ms & 2\,ms & 0.13\,s & 80{,}384\,B  & 260\,B & 128\,B \\
$2^{11}$ & 158.39\,s & 9.94\,s  & 1.51\,s  & 7.83\,s & 0.54\,s & 1\,ms & 17.16\,s & 0.43\,s & 2\,ms & 3\,ms & 0.13\,s & 98{,}624\,B  & 260\,B & 128\,B \\
$2^{12}$ & 316.07\,s & 19.61\,s & 5.13\,s  & 27.17\,s & 2.68\,s & 1\,ms & 34.15\,s & 0.81\,s & 2\,ms & 3\,ms & 0.13\,s & 119{,}168\,B & 260\,B & 128\,B \\
$2^{13}$ & 637.66\,s & 39.42\,s & 18.09\,s & 99.00\,s & 8.37\,s & 1\,ms & 66.95\,s & 1.49\,s & 2\,ms & 4\,ms & 0.13\,s & 135{,}872\,B & 260\,B & 128\,B \\
\bottomrule
\end{tabular}
\end{adjustbox}
\vspace{6pt}
\caption{Component-level measurements for \(\protocol\).}
\label{tab:lamp-breakdown}
\end{table*}

\subsection[LAMP Component Breakdown]
  {\rev{LAMP Component Breakdown}}
\label{sec:eval-lamp-breakdown}

\begin{revblock}
Table~\ref{tab:lamp-breakdown} presents a component-wise breakdown of the
costs of \(\protocol\) with the code rate set to \(\rho=1/2\). For small
matrices, Groth16-based proof generation is the dominant cost. As \(k\)
increases, however, encoding and commitment generation account for a larger
fraction of the total cost. At \(k=2^8\), encoding and commitment generation
take \(0.04\) s and \(0.23\) s, respectively, whereas generating the Merkle,
Groth16, and CP-Link proofs takes \(2.67\) s in total. In contrast, at
\(k=2^{13}\), encoding and commitment generation take \(125.46\) s in total,
exceeding the proof-generation cost of \(68.44\) s. This result shows that the
primary bottleneck shifts from proof generation to encoding and commitment
generation as the matrix dimension increases.

\(\protocol\) requires trusted setup for both Groth16 and CP-Link, with the
setup cost dominated by Groth16. The Groth16 setup time increases from
\(20.03\) s at \(k=2^8\) to \(637.66\) s at \(k=2^{13}\), whereas the CP-Link
setup time increases from \(1.28\) s to \(39.42\) s over the same range.

Verification time remains nearly constant as the matrix dimension increases.
Groth16 and Merkle verification each take only a few milliseconds, while
CP-Link verification takes approximately \(0.13\)--\(0.14\) s. The Groth16
and CP-Link proof sizes also remain constant at \(260\) B and \(128\) B,
respectively. Consequently, the total proof size is dominated by the Merkle
proof, which increases from \(49{,}920\) B to \(135{,}872\) B over the
measured range.

Table~\ref{tab:lamp-code-rates} in
Appendix~\ref{app:lamp-code-rates} presents additional results for
\(\rho\in\{1/2,1/4,1/8\}\). A higher code rate reduces encoding and commitment
costs by shortening the codeword, but it also decreases the relative distance
and therefore requires more queries. Specifically, the required numbers of
queries for \(\rho=1/8\), \(1/4\), and \(1/2\) are \(t=155\), \(189\), and
\(309\), respectively. Thus, the code rate determines the trade-off between
encoding and commitment costs and the generation time and size of the sampled
proof.
\end{revblock}

\begin{table*}[!t]
\centering
  \ifshowrevisions
    \rowcolors{1}{yellow!20}{yellow!20}
  \fi
\scriptsize
\setlength{\tabcolsep}{2.4pt}
\begin{adjustbox}{max width=\textwidth}
\begin{tabular}{c c ccc ccc ccc}
\toprule
\multirow{2}{*}{\textbf{Claims}}
& \multirow{2}{*}{\textbf{Constraints}}
& \multicolumn{3}{c}{\textbf{Commit Time}}
& \multicolumn{3}{c}{\textbf{Prove Time}}
& \multicolumn{3}{c}{\textbf{Proof Size}} \\
\cmidrule(lr){3-5}\cmidrule(lr){6-8}\cmidrule(lr){9-11}
& & Encode & ABC & XYZ
& Merkle & Groth16 & CP-Link
& Merkle & Groth16 & CP-Link \\
\midrule
1  & 329{,}377   & 0.11\,s & 0.22\,s & 0.02\,s & 0\,ms & 2.62\,s  & 0.07\,s & 49{,}728\,B & & \\
2  & 649{,}929   & 0.11\,s & 0.38\,s & 0.02\,s & 0\,ms & 4.81\,s  & 0.12\,s & 50{,}240\,B & & \\
3  & 970{,}481   & 0.12\,s & 0.54\,s & 0.03\,s & 0\,ms & 6.49\,s  & 0.15\,s & 49{,}088\,B & & \\
4  & 1{,}291{,}033 & 0.12\,s & 0.69\,s & 0.04\,s & 0\,ms & 8.71\,s  & 0.18\,s & 49{,}984\,B & & \\
5  & 1{,}611{,}585 & 0.13\,s & 0.83\,s & 0.05\,s & 0\,ms & 10.49\,s & 0.22\,s & 50{,}304\,B & & \\
6  & 1{,}932{,}137 & 0.13\,s & 0.80\,s & 0.05\,s & 0\,ms & 12.43\,s & 0.26\,s & 49{,}600\,B & & \\
7  & 2{,}252{,}689 & 0.13\,s & 0.94\,s & 0.05\,s & 0\,ms & 14.92\,s & 0.30\,s & 49{,}600\,B & & \\
8  & 2{,}573{,}241 & 0.14\,s & 1.02\,s & 0.06\,s & 0\,ms & 16.97\,s & 0.45\,s & 49{,}792\,B & & \\
9  & 2{,}893{,}793 & 0.14\,s & 1.19\,s & 0.06\,s & 0\,ms & 18.83\,s & 0.47\,s & 49{,}600\,B & & \\
10 & 3{,}214{,}345 & 0.15\,s & 1.25\,s & 0.07\,s & 0\,ms & 20.61\,s & 0.56\,s & 49{,}856\,B & \multirow{-10}{*}{260\,B} & \multirow{-10}{*}{128\,B} \\
\bottomrule
\end{tabular}
\end{adjustbox}
\vspace{6pt}
\caption{Batching \(q\) independent matrix-multiplication claims
\(\{\mat{A}_i\mat{B}_i=\mat{C}_i\}_{i\in[q]}\) at fixed \(k=2^8\).}
\label{tab:lamp-batch}
\end{table*}

\subsection[Batching Matrix-Multiplication Claims]
  {\rev{Batching Matrix-Multiplication Claims}}
\label{sec:eval-batch}
\begin{revblock}
We evaluate batching multiple independent claims
\(\mat{A}_i\mat{B}_i=\mat{C}_i\). At each codeword position, the prover packs
the symbols of all claims into a single Pedersen leaf and constructs one shared
Merkle tree. Consequently, the number of Merkle membership proofs is independent
of the batch size.

Table~\ref{tab:lamp-batch} reports the component costs at \(k=2^8\) and
\(\rho=1/2\). As the number of claims increases, the constraint count and
Groth16 proving time grow, while the Merkle proof size remains approximately
\(50\) KB and verification time remains approximately \(0.13\) s. This is
because the sampled positions and the shared Merkle openings are independent of
the number of claims; minor size variations result from shared authentication
nodes in the multi-opening implementation.
\end{revblock}

\subsection[GPT-2 Workload]
  {\rev{GPT-2 Workload}}
\label{sec:eval-gpt2}

\begin{revblock}
Table~\ref{tab:lamp-gpt2} reports the results for a GPT-2 layer with sequence
length \(2^{10}\) and 36 matrix-multiplication claims. To accommodate the
dimensions of the matrices in the layer, \(\protocol\) uses \(\rho=1/4\) and
\(t=189\). As in Section~\ref{sec:eval-matmul-comparison}, proving times include
encoding and commitment generation, and proof sizes include commitment data.
For DualMatrix, we configure the inputs of its public implementation to match
the GPT-2 matrix structures and evaluate the same matrix multiplications.

Compared with the Freivalds baseline, \(\protocol\) reduces the constraint
count from \(76{,}807{,}671\) to \(46{,}730{,}248\) and the proving time from
\(408.69\) s to \(367.40\) s. DualMatrix requires \(519.97\) s, making
\(\protocol\) the fastest prover among the three systems. In contrast, the
Freivalds baseline provides the fastest verification and smallest proof, while
DualMatrix also verifies faster and produces a smaller proof than \(\protocol\).

The relatively large proof size of \(\protocol\) in the GPT-2 experiment stems
from the different dimensions of the matrices involved in the workload. These
matrix multiplications require multiple commitment layouts and corresponding
CP-Link proofs and Merkle membership proofs. These linking and membership
proofs---dominated in size by the Merkle openings---therefore account for a
substantial portion of the total proof and are the main source of proof-size
growth in heterogeneous workloads such as GPT-2.
\end{revblock}

\begin{table}[H]
\centering
\ifshowrevisions
  \rowcolors{1}{yellow!20}{yellow!20}
\fi
% \footnotesize
\setlength{\tabcolsep}{3pt}
\begin{tabular*}{\columnwidth}{@{\extracolsep{\fill}}ccccc@{}}
\toprule
\textbf{System} & \textbf{Constraints} & \textbf{Prove} & \textbf{Verify} & \textbf{Proof Size} \\
\midrule
\(\protocol\) & 46{,}730{,}248 & 367.40\,s & 10.53\,s & 6{,}298{,}468\,B \\
Freivalds & 76{,}807{,}671 & 408.69\,s & 0.49\,s & 196\,B \\
DualMatrix & -- & 519.97\,s & 5.22\,s & 2{,}934{,}952\,B \\
\bottomrule
\end{tabular*}
\vspace{6pt}
\caption{GPT-2 at sequence length \(2^{10}\) with 36 matrix-multiplication
claims.}
\label{tab:lamp-gpt2}
\end{table}
